\newpage
\section{Wnioski}

Trzy funkcje zeszły na sterownik tą samą drogą. Tablica prawdy daje kanoniczną sumę, minimalizacja skraca ją do postaci z polecenia, a drabinka jest już rysunkiem tej postaci: szereg to iloczyn, węzeł to suma, kreska na styku to negacja. Po sklejeniu tablica się nie zmienia. Krótszy zapis ma tyle samo jedynek, tylko mniej styków.

Warto trzymać wspólny czynnik za węzłem, a nie przepisywać go w każdej gałęzi. W zadaniu~2 jedno $\overline{C}$ odcina i $A$, i $B$. W zadaniu~3 jedno $C$ puszcza dalej zarówno parę progów, jak i start ręczny. Powtórzenie tego styku w obu gałęziach dałoby tę samą tablicę i dłuższy program.

Ten sam symbol w zapisie algebraicznym nie mówi jeszcze, jakim stykiem wejdzie na drabinkę. W zadaniu~1 sygnał $C$ jest osobną gałęzią i sam załącza wyjście. W zadaniu~2 stoi w szeregu jako styk zanegowany i wyjście gasi. W zadaniu~3 też jest w szeregu, ale stykiem zwykłym, bo zasilanie ma być obecne. O tym, która wersja jest właściwa, decyduje tablica, a nie podobieństwo litery.

Symulator wystarczył, żeby zobaczyć, która gałąź trzyma cewkę: sam iloczyn, sama suma, sama blokada. Żarówka przy Q1 zgadza się z wierszem tablicy dla ustawionej kombinacji. Nie zastępuje to podłączonego przycisku ani styku, który na obiekcie bywa odwrotnie spolaryzowany. Tę część instrukcja zostawia na test sprzętowy, a w tym sprawozdaniu jej nie ma.
